% XeLaTeX chapter fragment. Dependencies: ctexbook/ctexrep, amsmath, amssymb,
% booktabs, graphicx. Load generated_tables.tex once before these chapters.
\chapter{基于执行状态与恢复后缀配对监督的小模型知识图谱问答方法}
\label{chap:execution-feedback}

\section{研究问题与方法定位}

知识图谱问答需要将自然语言问题转换为可执行的查询操作。分步执行的模型能够观察中间实体集合、属性值与执行失败信息，但这并不意味着模型已经学会利用这些状态继续求解。仅以正确参考轨迹进行监督训练时，训练上下文主要由参考动作及其执行结果组成；自由生成时，模型可能进入训练中较少出现的状态，并在此后反复调用无效操作或耗尽调用预算。本章研究的问题是：在相同初始模型、训练问题来源和监督 token 预算下，为小模型补充真实自由运行状态及其可执行后续轨迹，能否提高问答成功率？

本文构建一种恢复轨迹监督训练方案。先从初始模型在训练集上的失败轨迹中，定位其与参考轨迹的首次分歧；随后保留此前成功执行的共同前缀，将真实分歧动作及其反馈作为上下文，重放经过句柄重映射的参考后缀，并仅对后缀中的模型输出设置监督目标。与之配对的普通续训样本使用相同问题、相同共同前缀及相同参考后缀，但不插入该分歧动作。两组再混入完全相同的普通成功轨迹，并匹配逐记录监督 token 数、记录次序和重复次数。

本章的方法贡献位于\textbf{恢复状态选择、可执行后缀构造与监督位置控制}，使用已有的语言模型、LoRA 和交叉熵训练，不声称提出新的网络结构或优化器。恢复训练的整体效果与具体反馈机制分别评价：即使整体方法改善最终准确率，也不能直接推导出模型理解了报错文字、识别了语义错误或完整保留了问题中的所有条件。

\section{执行环境与任务形式化}

设知识图谱为 $G$，自然语言问题为 $q$，执行状态为 $s_t$。状态包含当前已执行的动作、环境返回值及可引用的中间结果句柄。模型根据问题和历史 $h_t$ 生成动作 $a_t$，执行器返回观察 $o_t$ 并更新状态：
\begin{align}
  a_t &\sim \pi_\theta(\,\cdot\mid q,h_t), \\
  (s_{t+1},o_t) &= \mathcal{E}(G,s_t,a_t), \\
  h_{t+1} &= h_t\mathbin{\Vert}(a_t,o_t).
\end{align}
其中 $\mathcal{E}$ 是固定执行器，$\Vert$ 表示追加历史。动作包括知识操作及提交答案的结束操作；中间结果通过句柄引用，后续操作可以使用已有结果作为依赖。执行器返回成功标志、结果内容或失败说明。某个操作能够合法执行，并不保证它符合问题意图；空结果也可能是合法反馈。

自由运行产生轨迹 $\tau=(a_1,o_1,\ldots,a_T,o_T)$。如果模型在固定预算内提交了答案，则根据公开任务所用答案比较函数判断最终答案是否正确。预算耗尽、执行异常或未能提交答案的题目均保留在评价分母中。训练样本的正确性核验与测试时的自由生成严格分离：恢复样本构造可使用训练题的参考程序和答案，而开发集及留出集推理只向模型提供问题及允许的执行接口。

\section{真实失败轨迹的配对构造}

\subsection{共同前缀与分歧位置}

对训练问题 $q$，设参考操作序列为 $a^{*}_{1:K}$，初始模型自由生成的序列为 $\hat a_{1:T}$。首先回放自由运行轨迹，检查实际执行事件与记录一致，并确认该题最终未正确作答。随后寻找从起点开始、与参考动作及成功观察匹配的最长共同前缀，长度记为 $k$。构造要求 $k\geq 1$，并要求存在共同前缀之后的首次分歧动作 $\hat a_{k+1}$。

共同前缀长度采用参考轨迹匹配确定，因此首次分歧是\textbf{样本选择标记}，不是已标注的语义错误位置。参考程序并非唯一可能的正确程序；一个与参考不同的动作仍可能是合理的探索或另一条求解路径。本文只使用最终失败且满足构造约束的轨迹，不把所有分歧动作标为错误。

\subsection{普通后缀与恢复后缀}

将共同前缀对应的历史记为 $H_k$。普通续训组 A 和恢复续训组 C 的配对序列写为
\begin{align}
  x^{A} &= (q,H_k,a^{*}_{k+1:K},\mathrm{Done}), \\
  x^{C} &= (q,H_k,\hat a_{k+1},\hat o_{k+1},
             \rho(a^{*}_{k+1:K}),\mathrm{Done}).
\end{align}
式中动作之间均包含真实执行器返回的观察，为简洁未在每一项展开；$\rho$ 为中间结果句柄重映射。分歧动作可能产生额外句柄，使得后续结果编号改变。因此，C 中参考后缀的函数、常量参数及依赖图保持不变，但依赖句柄按当前执行状态重新映射，不能简单复制 A 的文本序列。

构造器分别从干净前缀和含分歧动作的前缀继续执行后缀。只有在后缀可执行、执行结果与对应参考结果匹配、最终答案通过比较时，才保留该配对。若插入的分歧动作已经终止回合、超过调用预算、无法匹配前缀或不能重放参考后缀，则排除该样本。此检查证明的是构造出的训练目标可在规定环境中执行，不能替代人工对自然语言条件的语义核验。

\subsection{构造流程与实际覆盖}

构造过程可概括为以下步骤。
\begin{enumerate}
  \item 固定训练题 ID、执行环境、参考程序和初始模型，在训练题上进行一次自由运行。
  \item 回放轨迹并核对事件，从最终失败的轨迹中识别至少一个成功参考动作组成的共同前缀。
  \item 提取首次参考分歧动作及其实际观察，分别构造 A 的普通后缀与 C 的恢复后缀。
  \item 执行句柄重映射，完整重放两条目标轨迹，核验后缀结果及最终答案。
  \item 记录样本来源、排除原因和监督掩码，冻结训练配对与来源摘要。
\end{enumerate}

在 $5\,000$ 条训练自由运行轨迹中，$4\,998$ 条通过记录回放核验，其中 $659$ 条最终失败、$4\,339$ 条正确。最终保留 $418$ 个真实恢复配对，覆盖全部训练问题的 $8.36\%$。主要排除项包括 $240$ 条缺少共同成功前缀、$2$ 条回放记录不一致及 $1$ 条参考答案核验不通过。$418$ 个配对中，分歧动作有 $414$ 个可以执行、仅 $4$ 个被执行器拒绝。因此，这批数据主要是\textbf{可执行分歧状态之后的恢复监督}，不能描述成以大量语法报错修复样本为主体的数据集。

上述计数来自冻结的配对构造报告；合格样本较少是当前方法的覆盖限制。为达到训练预算而重复使用同一配对，不会增加独立问题数量。未来领域迁移需要重新检查数据覆盖，不能把公开数据上的比例直接外推到雷达数据。

\section{监督掩码与匹配预算}

\subsection{动作后缀监督}

对序列 $x=(x_1,\ldots,x_L)$，令 $m_j\in\{0,1\}$ 为监督掩码。共同前缀、真实分歧动作、所有工具观察以及提示文本的掩码均为零；后续参考动作及结束输出的模型 token 初始设为一。训练使用因果语言建模损失：
\begin{equation}
  \mathcal{L}_{\mathcal B}(\theta)
   =-\frac{1}{\sum_{x\in\mathcal B}\sum_j m_j}
     \sum_{x\in\mathcal B}\sum_j
       m_j\log p_\theta(x_j\mid x_{<j}),
\end{equation}
其中 $\mathcal B$ 表示计入相应优化步骤的训练记录，实际归一化和梯度累计由固定版本的 Trainer 实现。此式说明监督对象；它不意味着整个训练过程等价于一次性对全语料求解全局平均损失。

分歧动作仍然出现在 C 的输入中，模型可读取其内容及结果，但不被要求模仿它。工具观察也作为条件输入而非生成目标。普通轨迹部分在 A/C 中完全一致；配对后缀部分的主要区别为 C 所见的额外执行状态及相应句柄引用。

\subsection{逐记录 token 匹配}

为避免 C 仅因获得更多目标动作而占优，使用训练成功轨迹一个遍历所含的 $788\,534$ 个监督 token 作为续训总预算。预算一半用于普通成功轨迹，一半用于配对后缀，各为 $394\,267$。对每个配对记录，目标上限为
\begin{equation}
  b_i=\min\{n_i^A,n_i^C,B_{\mathrm{remain}}\},
\end{equation}
其中 $n_i^A,n_i^C$ 是两组该记录原有的监督 token 数，$B_{\mathrm{remain}}$ 是该部分剩余预算。超出上限的目标按时间顺序从尾部屏蔽，输入上下文保持完整。两组共享问题、记录次序、重复次数及逐记录监督数量，使用固定语料构造种子。

\begin{table}[htbp]
  \centering
  \caption{A/C 冻结续训语料的计数。输入 token 为语料输入长度总和。}
  \label{tab:matched-corpus}
  \PublicTrainingTable
\end{table}

两组各有 $5\,772$ 条训练记录，覆盖 $2\,694$ 个不同问题。普通部分使用 $2\,486$ 条记录；配对部分为 $3\,286$ 条记录，源于 $418$ 个不同配对，其中 $58$ 个各出现七次、$360$ 个各出现八次。这些重复记录是预算安排，不能按独立样本计数。

预算追加屏蔽的 token 在 A/C 中分别为 $194$ 和 $321$。C 额外的 $126$ 次小幅尾部屏蔽对应换行符；少数较大裁剪仍可能屏蔽动作末尾的部分目标。故不能声称每条记录中的目标动作始终保持完整。两组输入 token 总数分别为 $9\,135\,350$ 和 $9\,520\,876$，可见\textbf{相同监督 token 预算不等于相同上下文长度、浮点计算量或训练时间}。

\section{训练与推理实现}

模型使用 Qwen2.5-3B-Instruct。先以通过回放核验的成功轨迹训练初始分步模型 I，再从同一个 I 检查点分别续训 A 和 C。LoRA 秩为 $16$，缩放参数为 $32$，dropout 为 $0.05$，作用于实现中指定的全部线性层；使用现有 AdamW 优化器。A/C 学习率均为 $10^{-4}$，每设备批大小为 $2$，梯度累计步数为 $8$，在冻结的续训语料上训练一个遍历。上下文上限为 $8\,192$，使用 BF16 与梯度检查点。参数及依赖版本以训练输入清单和源码摘要为准。

两组分别使用续训随机种子 $20261003$ 和 $20261004$ 进行配对实验。两次均继承相同的初始模型与训练语料；它们评价的是改变\textbf{续训随机性}后的效果，不构成两次从基础模型开始的全流程独立复现。语料构造种子、推理种子及 Python 哈希种子固定为 $20261003$。

推理使用贪心解码，批大小为 $8$，每题最多 $24$ 次工具调用，总生成上限为 $2\,048$ 个 token，单轮生成上限为 $192$，上下文上限为 $8\,192$。模型读取原始执行反馈，不在留出评测阶段修改提示、工具接口或反馈文本。成本指标累计每轮完整输入和生成 token，生成部分包含 EOS；输入中重复出现的历史也重复计入。

\section{可检验贡献与适用边界}

本方法提出可复现的配对训练构造，并把两类问题分开检验。第一类是方法整体是否有效，即在匹配监督预算后 C 是否优于 A；第二类是收益如何产生，例如是否减少未完成回合、是否依赖详细报错文字。整体效果由独立项目留出集评价，机制解释由开发集对照辅助分析。后者证据不足时，仍须保留不确定结论。

当前证据所能支持的贡献是：\textbf{通过真实分歧状态与可执行参考后缀的配对构造，将训练监督集中于恢复后的模型输出，并在匹配监督 token 的对照下检验其问答收益。}这是一项训练数据与监督方式的研究改进。状态插入、句柄变化和恢复后缀监督作为整体方案被比较，尚未通过独立消融分离每个组成部分的贡献；也未证明比所有现有恢复训练方法更优。

该方法适用于具有确定执行接口、训练参考轨迹和可回放中间结果的知识图谱问答任务。迁移到雷达领域时，还需处理型号别名、指标单位、工作条件与来源引用。公开实验不能代替雷达事实核验，也不能直接作为雷达领域效果。后续领域实验应固定可信事实及独立问题，再检验相同训练思想在这些约束下是否仍有收益。

\section{本章材料的证据入口}

配对数量与选择规则来自冻结的 \texttt{recovery\_pairs.json}；语料与监督预算来自 \texttt{continuation\_data.json} 及 \texttt{continuation\_token\_audit.json}；模型设置来自各训练运行的 \texttt{input\_manifest.json} 与实现。具体文件路径、SHA-256 和表格导出规则记录在本目录的 \texttt{evidence\_manifest.json} 与 \texttt{results\_tables.csv} 中。本文没有在写作阶段新增训练样本、重新选择模型或打开逐题留出预测。
